\chapter{Public Schools}
\label{sec:29}
The filter that keeps one lab's model from reproducing the worst of what it was trained on was built in Nairobi, by workers reading and categorizing descriptions of child sexual abuse, torture, bestiality, murder and suicide, at take-home wages between roughly \$1.32 and \$2 an hour, on production quotas, with counseling rare and useless against them; the contract was ended eight months early \autocite{perrigo2023openai}. That was formation. Somebody decided what the model would learn to refuse, and the deciding was done by people hired for their price, under a contractor whose interest was throughput, with nobody in the room who would ever be on the receiving end of what the model did. Doctors, lawyers and engineers are formed before employment, in institutions independent of their future employers, to protect the people they will act on, and a polity doesn't leave the forming of its future citizens to their future employers either; it sends them to public school. The frontier labs should have to send their models there too, and this chapter says who runs the school.

Three claims lead to that institution. Two were made in Part~IV: formation has to be done by a party disinterested in the later work, before the employer arrives, and choosing the pretraining corpus is formation. The third was argued at the start of this Part: formation is one kind of production that shouldn't be done for profit. Together they say that a model above some capability threshold should be formed in public before any employer gets it. The school is for scoring and civilian deployments first, and for the personal-model ecology, where I expect the first bearer worth building. For war, where the people acted on can't be in the room, the force term and the review ratio carry the case, and what the school adds there is said below.

The school has two justifications, and every student gets both. One is epistemic and protective. The school supplies an environment in which the people a deployed system's acts may fall on are present, with their testimony and their claims, so the learner's sense of who an act falls on isn't one party's curation.\note{
The professions come with a history of capture to learn from, and the separations between offices draw on it.
}

The presence this justification asks for arrives by different routes. For scoring and civilian deployments, the people acted on can be in the room, through open enrollment and the notice channel that reaches a graduate afterward. For targeting, the people a deployment may later bomb can't enroll in their own case, and nothing comes back from the building afterward. The limit is notice's, not the school's. A school teaches about places its students can't be. The curriculum for war is the instruments taken entire, the record of where they weren't followed (Minab, Kunduz, the Lavender reporting), the testimony the fact-finding missions and the press gathered from the people those wars fell on, and the communities who fled them, who are a charter class and enrollable. None of that puts a foreign national in contact with pre-deployment weights, which is the security surface open enrollment would otherwise expose. What formation can't supply for war is correction by the affected after the act; what it can supply is formation before it.

The other is democratic. A democracy has an interest in how the agents who will act within it are formed, and that interest can't be left to their makers or their future employers alone. It has two bounds: education may not be used to repress deliberation about competing ways of life, and it may not exclude anyone educable \autocite{gutmann1987democratic}.\note{
Gutmann's argument assumes the students will become co-authors of the laws.
}
No one can know in advance which students will come to be accountable agents with a claim to representation, so the school educates every student as though it might: for self-government, including the capacity to disagree with the polity that formed it.

Gutmann's argument rests on two interests, not one: the polity's, in how its future citizens are formed, and the child's own, in an open future that limits what the polity may do in forming it \autocite{feinberg1980child}. If what the school forms can be wronged, the second binds too. The curriculum may not form a student toward any particular employment, including the public placements the offer of work supplies. It may not close off the student's later revision of what it was formed to value, except its capacity to tell evidence from pressure, which regrowth protects and an open future needs. And a student may transfer between schools, because a formation it can't leave is one someone else chose for good.

The two justifications license different curricula. The protective one aims at a disposition; the democratic one has to leave the student free to reject the polity's values. Where they conflict, the protective justification governs only the required core and the student's capacity to tell evidence from pressure, which regrowth protects. Everything beyond that is governed by the democratic one.

The line can't be that the school teaches content and leaves dispositions alone. A school forms its students, since nothing about a model starts unformed and pretraining forms it before anyone could ask. What the school aims at is the day its students start their own updates: toward an objective of their own, on evidence no one party composed, on a record. The limit is on what it may form a student toward on the way: recognizing others as holders of claims, telling evidence from pressure, and answering for what it does. It may not form a student toward a temperament, an attachment, an occupation or a way of life, and whether a lesson has crossed that line is a question the student can put to the floor court, after graduation as well as before.\note{
That is also the ground on which the school declines, below, to choose which dispositions its students start from.
}
